\begin{theorem}[Complete charge measurement on an actual adjacent pair]
    \label{thm:peps_regular_edge_physical_charge_measurement}
    \lean{TNLean.PEPS.exists_regularEdgePhysicalAllChargeMeasurement}
    \leanok
    \uses{thm:peps_regular_edge_charge_contraction,
        thm:peps_regular_two_site_all_charge_measurement}
    Assume the two original site maps at the endpoints of $e$ are regular
    $G$-isometric. Let $\mathcal C$ be the actual irreducible characters of the
    regular representation. There is a complete physical projective measurement
    $(Q_r)_{r\in\mathcal C\cup\{\varnothing\}}$ on these two sites, with
    $Q_r^*=Q_r\geq0$, $Q_rQ_s=\delta_{r,s}Q_r$, and $\sum_rQ_r=I$, such that
    \begin{align}
        Q_r\Psi^{\chi,p}_{R_e,\theta}
        =\mathbf1_{r=\chi}\Psi^{\chi,p}_{R_e,\theta}.
    \end{align}
    The same family is chosen before every $\chi\in\mathcal C$, $p$, $\theta$,
    and every exterior site family whose tensors agree with the original ones at
    both measured endpoints. The outcome $\varnothing$ annihilates every charged
    column. This is a finite-graph form of
    \cite[Detection of chargeons]{Schuch2010PEPS}, lines 2464--2486.
\end{theorem}
\begin{proof}\leanok
    Transport the complete two-site physical measurement along the actual
    endpoint-coordinate bijection. The literal charged coefficient identity
    identifies all native columns with the two-site charged columns. The
    bijection preserves adjoints, positivity, products, and the sum of projections.
    Only the two endpoint tensors enter this calculation.
\end{proof}
